\chapter{Lineær algebra}
Lineær algebra beskriver vektorrum og lineære afbildninger mellem dem. Dens styrke er, at mange problemer inden for mekanik, elektromagnetisme, differentialligninger og kvantemekanik kan samles i den samme struktur: vektorer, operatorer og deres invariant definerede egenskaber. En matrix er en \emph{repræsentation} af en lineær afbildning i valgte baser; afbildningen selv afhænger ikke af basisvalget.

\section{Vektorrum, basis og lineære afbildninger}
Et vektorrum $V$ over et legeme $\mathbb K$ (typisk $\mathbb R$ eller $\mathbb C$) er lukket under addition og skalarmultiplikation og opfylder vektorrumsaksiomerne. Et sæt $(v_1,\ldots,v_n)$ er en basis, hvis enhver $v\in V$ har en \emph{entydig} fremstilling
\[
v=\sum_{i=1}^{n}v^i v_i.
\]
Koefficienterne $v^i$ er koordinater i den valgte basis. De er ikke selve vektoren. I et endeligdimensionalt rum har alle baser samme antal elementer, nemlig $\dim V$.

En afbildning $T:V\to W$ er lineær, når
\[
T(av+bw)=aT(v)+bT(w)
\qquad(a,b\in\mathbb K).
\]
I baserne $(v_j)$ og $(w_i)$ defineres matricen $A=(A^i{}_j)$ ved
\[
T(v_j)=\sum_i A^i{}_j w_i,\qquad [T(v)]_w=A[v]_v.
\]
Ved et basisskift $[v]_{\mathrm{gammel}}=S[v]_{\mathrm{ny}}$ for en operator på samme rum gælder
\[
A_{\mathrm{ny}}=S^{-1}A_{\mathrm{gammel}}S.
\]
Dette er en \emph{similaritetstransformation}. Den bevarer karakteristisk polynomium, egenværdier, determinant og spor. Den må ikke forveksles med omregning af en bilineær form, hvor transformationsloven indeholder en transponeret matrix.

\subsection{Kerne, billede og dimension}
Kernen og billedet defineres ved
\[
\ker T=\{v\in V:T(v)=0\},\qquad
\operatorname{im}T=\{T(v):v\in V\}.
\]
For endeligdimensionalt $V$ giver dimensionssætningen
\[
\boxed{\dim V=\dim(\ker T)+\operatorname{rank}T.}
\]
En lineær afbildning er injektiv præcis når dens kerne kun består af nulvektoren. For en kvadratisk matrix er injektivitet, surjektivitet og invertibilitet ækvivalente.

\section{Determinanter, rang og lineære ligninger}
Et lineært ligningssystem $Ax=b$ er konsistent, hvis og kun hvis $b$ ligger i søjlerummet for $A$, eller ækvivalent
\[
\operatorname{rank}(A)=\operatorname{rank}([A\mid b]).
\]
Hvis løsninger findes, er deres samlede form $x=x_{\mathrm p}+z$, hvor $Az=0$. Antallet af frie parametre er $\dim\ker A$. Gauss-elimination ændrer ikke løsningsmængden, når man udfører elementære rækkeoperationer på den udvidede matrix.

Determinanten findes for kvadratiske matricer. Dens absolutværdi er volumenfaktoren for den tilsvarende lineære transformation, mens fortegnet i reelle rum angiver orienteringsskift:
\[
\det(AB)=\det A\det B,\qquad
\det(A^{\mathsf T})=\det A,\qquad
\det(A^{-1})=\frac{1}{\det A}\quad(\det A\neq0).
\]
En kvadratisk matrix er invertibel præcis når dens determinant er forskellig fra nul. Determinanten kan dog være numerisk uhensigtsmæssig som praktisk test for næsten-singularitet: konditionstal og singularværdier giver ofte en bedre vurdering.

\subsection{Spor og karakteristisk polynomium}
Sporet er summen af diagonalelementerne og er basisuafhængigt:
\[
\operatorname{tr}(A+B)=\operatorname{tr}A+\operatorname{tr}B,\qquad
\operatorname{tr}(AB)=\operatorname{tr}(BA).
\]
For en $n\times n$-matrix er det karakteristiske polynomium
\[
p_A(\lambda)=\det(\lambda I-A).
\]
Over $\mathbb C$ er summen af egenværdierne med algebraisk multiplicitet $\operatorname{tr}A$, og produktet er $\det A$.


\section{Fra diagonalelementer til spor og tensorkontraktion}
\label{sec:spor-kontraktion}
En vigtig idé, som vender tilbage i tensorregning og differentialgeometri, er at danne en skalar ved at summere over et par sammenhørende indekser. I lineær algebra møder vi den først som \emph{sporet} af en lineær operator.

For en kvadratisk matrix $A=(A^i{}_j)$ er
\[
\boxed{\operatorname{tr}A=\sum_{i=1}^{n}A^i{}_i.}
\]
Summen løber over \emph{diagonalelementerne}, men sporet er ikke det samme som at finde eller konstruere en diagonal matrix. Som en første illustration er
\[
A=\begin{pmatrix}2&5\\-1&3\end{pmatrix},
\qquad
\operatorname{tr}A=2+3=5.
\]
Derimod betyder \emph{diagonalisering} at finde en basis, hvor matrixrepræsentationen kan skrives $D=P^{-1}AP$ med alle elementer uden for diagonalen lig nul. Dette kræver tilstrækkeligt mange lineært uafhængige egenvektorer.

Hvorfor er sporet uafhængigt af basis? Under et basisskift bliver
$A'=S^{-1}AS$. Da $\operatorname{tr}(BC)=\operatorname{tr}(CB)$ for kvadratiske matricer med kompatible dimensioner, får vi
\[
\operatorname{tr}(S^{-1}AS)
=\operatorname{tr}(ASS^{-1})
=\operatorname{tr}A.
\]
Den specifikke diagonal ændrer sig altså ved basisskift, men \emph{summen} af dens elementer gør ikke.

I tensorregning generaliseres dette til \emph{kontraktion} af et øvre og et nedre indeks, f.eks.
\[
T^i{}_i=\sum_i T^i{}_i.
\]
I dette eksempel er kontraktionen netop sporet. For en tensor med flere indekser kan én kontraktion efterlade en tensor af lavere orden. Man kan ikke generelt kontrahere to vilkårlige øvre indekser uden yderligere struktur: en metrik kan bruges til at hæve eller sænke indeks, så parringen giver geometrisk mening.

\paragraph{Forbindelse fremad.}
I bind 5, kapitlet om tensorregning, bruges samme idé til f.eks.
$R=g^{ij}R_{ij}$: Ricciskalaren dannes ved at kontrahere Riccitensoren med den inverse metrik. Sporet af en matrix er derfor et første, konkret tilfælde af en langt bredere geometrisk operation.

\paragraph{Hvornår bruger man hvad?}
Vil man beregne summen af egenværdierne med algebraisk multiplicitet, bruger man sporet. Vil man finde de retninger, hvor en lineær operator virker ved simpel skalering, undersøger man egenvektorer og diagonalisering. Vil man reducere antallet af frie indekser i en tensorlig relation på en koordinatuafhængig måde, bruger man kontraktion.

\section{Egenværdier og diagonalisering}
En ikke-nul vektor $v$ kaldes egenvektor for $A$, hvis $Av=\lambda v$. Egenværdierne opfylder
\[
\boxed{\det(A-\lambda I)=0.}
\]
Egenvektorer med indbyrdes forskellige egenværdier er lineært uafhængige. Derfor kan en $n\times n$-matrix diagonaliseres, hvis den har $n$ lineært uafhængige egenvektorer:
\[
A=PDP^{-1},\qquad D=\operatorname{diag}(\lambda_1,\ldots,\lambda_n).
\]
Det er ikke tilstrækkeligt blot at have $n$ egenværdier \emph{talt med algebraisk multiplicitet}: en matrix kan mangle egenvektorer. For en sådan matrix kan Jordan-normalform være relevant.

Når diagonalisering findes, kan man beregne funktioner af matricen ved at anvende funktionen på diagonalelementerne:
\[
f(A)=P\,\operatorname{diag}\bigl(f(\lambda_1),\ldots,f(\lambda_n)\bigr)P^{-1}.
\]
For den lineære differentialligning $\dot x=Ax$ giver dette særligt
\[
x(t)=e^{At}x(0),\qquad
e^{At}=Pe^{Dt}P^{-1}.
\]
Matrixeksponentialen er også defineret, når $A$ ikke kan diagonaliseres, ved sin absolut konvergente potensrække.

\section{Indre produkter, spektralsætning og SVD}
Et indre produkt $\langle u,v\rangle$ giver et mål for længde og ortogonalitet. På $\mathbb R^n$ bruges typisk $u^{\mathsf T}v$ og på $\mathbb C^n$ det hermiteske produkt $u^\dagger v$ (konjugeret lineært i første argument). Normen er
\[
\|v\|=\sqrt{\langle v,v\rangle},\qquad
|\langle u,v\rangle|\leq \|u\|\|v\|.
\]
En ortonormal basis opfylder $\langle e_i,e_j\rangle=\delta_{ij}$. I en sådan basis er den adjungerede operators matrix $A^\dagger$.

\subsection{Selvadjungerede og normale operatorer}
En kompleks matrix er hermitesk, hvis $A^\dagger=A$, og en reel matrix er symmetrisk, hvis $A^{\mathsf T}=A$. Den endeligdimensionale spektralsætning siger, at en normal kompleks matrix ($AA^\dagger=A^\dagger A$) kan diagonaliseres ved en unitær matrix:
\[
A=UDU^\dagger,\qquad U^\dagger U=I.
\]
For en hermitesk matrix er egenværdierne reelle. Egenvektorer svarende til forskellige egenværdier er ortogonale; dette er centralt for blandt andet normale svingninger og kvantemekaniske observabler.

\subsection{Singularværdidekomposition}
Enhver matrix $A\in\mathbb C^{m\times n}$ har en singularværdidekomposition
\[
\boxed{A=U\Sigma V^\dagger,}
\]
hvor $U,V$ er unitære, og $\Sigma$ er rektangulær diagonal med ikke-negative singularværdier $\sigma_i$. De ikke-nul singularværdier er kvadratrødderne af de positive egenværdier for $A^\dagger A$. Antallet af positive singularværdier er matricens rang.

SVD findes også, når matricen ikke er kvadratisk eller diagonaliserbar. Den er derfor et generelt værktøj til projektion, rangvurdering og mindste-kvadraters problemer. For en matrix med fuld søjlerang er mindste-kvadraters løsningen til $\min_x\|Ax-b\|_2$ entydig; uden fuld rang vælges typisk den løsning, der har mindst norm, ved hjælp af pseudoinversen $A^+$.

\subsection{Ortogonale projektioner}
For en ortonormal søjlematrix $Q$ med $Q^\dagger Q=I$ er
\[
P=QQ^\dagger,\qquad P^2=P,\qquad P^\dagger=P.
\]
$P$ projicerer ortogonalt på søjlerummet for $Q$. Det gør relationen mellem lineær algebra, geometri og approksimation tydelig: en bedste approksimation i et underrum fås ved at fjerne den ortogonale restkomponent.

\paragraph{Anvendelser i fysik.}
Normalmoders frekvenser findes typisk fra generaliserede egenværdiproblemer; lineæriserede bevægelsesligningers stabilitet bestemmes blandt andet af operatorens spektrum; ortonormale baser og selvadjungerede operatorer er centrale i kvantemekanik. For ikke-ortogonale koordinater skal metrikken indgå eksplicit i indre produkter og adjungering.

\section{Praktisk metodevalg: lineære systemer og egenværdier}
\label{sec:la-metode}
Inden man regner, bør man skelne mellem tre forskellige spørgsmål. Et ligningssystem $Ax=b$ handler om en \emph{løsning}; egenværdiproblemet $Av=\lambda v$ handler om operatorens \emph{naturlige retninger}; og et mindste-kvadraters problem handler om den \emph{bedst mulige approksimation}, når ligningerne ikke kan opfyldes præcist.

\subsection{Sådan vælges metoden}
\begin{center}
\begin{tabular}{@{}p{0.32\textwidth}p{0.57\textwidth}@{}}
\toprule
Spørgsmål & Typisk fremgangsmåde\\
\midrule
Er $Ax=b$ løsbart? & Sammenlign $\operatorname{rank}A$ og $\operatorname{rank}[A\mid b]$.\\
Er løsningen entydig? & Kræv en konsistent ligning og $\ker A=\{0\}$.\\
Er en kvadratisk matrix invertibel? & Brug rang eller determinant; numerisk også singularværdier.\\
Hvilke naturlige moder har systemet? & Stil egenværdiproblemet op, kontroller massen og fortegnene.\\
Kan operatoren diagonaliseres? & Kontroller antal lineært uafhængige egenvektorer; ikke blot antal rødder.\\
Hvilken approximation er nærmest $b$? & Benyt ortogonal projektion eller SVD.\\
\bottomrule
\end{tabular}
\end{center}
I numeriske beregninger bør man normalt løse $Ax=b$ med en faktorisering eller eliminationsmetode i stedet for først at beregne $A^{-1}$. At inversen findes matematisk, siger ikke alene at problemet er numerisk stabilt.

\subsection{Matrix og koordinatvalg}
Lad $v$ være en geometrisk vektor. Den kan have forskellige koordinatsøjler $[v]_e$ og $[v]_f$ i to baser. Skriver vi $[v]_e=S[v]_f$, følger for en operator $T:V\to V$
\[
[T]_f=S^{-1}[T]_eS.
\]
Denne formel forklarer, hvorfor eigenværdier og spor er basisuafhængige, selv om alle de enkelte matrixelementer kan ændre sig. Hvis man i stedet skifter baser på både domæne og codomæne i en afbildning $T:V\to W$, bliver den generelle regel $[T]_{\mathrm{ny}}=R^{-1}[T]_{\mathrm{gammel}}S$, hvor $S$ og $R$ beskriver skiftene i henholdsvis $V$ og $W$.

\subsection{Kvadratiske former og positive matricer}
En kvadratisk form kan skrives
\[
Q(x)=x^{\mathsf T}Kx.
\]
Kun den symmetriske del af $K$ bidrager, da $x^{\mathsf T}(K-K^{\mathsf T})x=0$. For en reel symmetrisk matrix $K$ gælder: $Q(x)>0$ for alle $x\ne0$ præcis når alle egenværdier af $K$ er positive. Denne karakterisering viser, hvorfor egenværdier bruges til at afgøre, om et stationært punkt er et stabilt lokalt minimum i en lineær eller kvadratisk model.

\section{Fysisk sammenhæng: koblede små svingninger}
\label{sec:normalmoder}
For små forskydninger $q$ omkring en ligevægt får man ofte en Lagrangefunktion på formen
\[
L=\frac12\dot q^{\mathsf T}M\dot q-\frac12q^{\mathsf T}Kq,
\]
hvor $M$ er en symmetrisk positiv definit massematrix, og $K$ er en symmetrisk stivhedsmatrix (Hessianen af potentialet i ligevægten). Euler--Lagrange giver
\[
M\ddot q+Kq=0.
\]
Ved ansatsen $q(t)=a\cos(\omega t+\varphi)$ fremkommer
\[
(K-\omega^2M)a=0,\qquad
\boxed{\det(K-\omega^2M)=0.}
\]
Dette er et \emph{generaliseret egenværdiproblem}. Det reduceres kun direkte til $Ka=\omega^2a$, når $M=I$ i de valgte enheder og koordinater. Hvis $K$ er positiv semidefinit, er de fundne $\omega^2$ ikke-negative; en nulmode kan svare til en fri bevægelsesretning. En negativ egenværdi af det relevante massenormaliserede problem giver i stedet en eksponentielt voksende mode.

Når $M$ er positiv definit, kan man indføre de massenormaliserede koordinater med $u=M^{1/2}a$ og den symmetriske matrix
\[
B=M^{-1/2}KM^{-1/2},
\qquad Bu=\omega^2u.
\]
Så kan den almindelige spektralsætning anvendes på $B$. Dette er forbindelsen mellem lineær algebra, andenordens differentialligninger og stabilitet i analytisk mekanik.

\subsection{Hvorfor superposition virker}
Da bevægelsesligningen er lineær, er summen af løsninger også en løsning. Når massenormaliserede egenvektorer danner en basis, kan de almindelige initialbetingelser opdeles på naturlige moder; hver mode svinger med sin egen frekvens. Superposition gælder i \emph{den lineære model}; store udsving kan kræve højere ordens led i potentialet, hvor normalmoderne ikke længere udvikler sig uafhængigt.

\section{Forbindelser til senere bind}
De matematiske resultater her er forudsætninger for en række senere emner. Jacobimatricen i flervariabel analyse er den lokale lineære afbildning, der bedst approksimerer et differentiabelt koordinatskift. Fouriermetoder skriver funktioner i en basis af oscillationer; Sturm--Liouville-teori generaliserer egenværdiproblemer til differentialoperatorer. Tensorregning skelner mellem objekter og deres koordinatkomponenter og generaliserer sporet ved kontraktion.

Den røde tråd er derfor ikke, at nye emner er det samme som lineær algebra, men at de ofte \emph{bevarer dens grundstruktur} under mere generelle betingelser.
